\section{Introduction}
\label{sec:intro}

Large language models generate text left-to-right and, even when equipped with
chain-of-thought \citep{wei2022chain}, commit to a single reasoning trajectory
without the ability to retract it.
When a reasoning path reaches a dead end, the model cannot backtrack and
explore an alternative branch; it either continues down an incorrect path or
abruptly changes direction without any explicit representation of the abandoned
subgoal.
This limitation is not merely a failure of prompting style---it reflects the
absence of any syntactic mechanism for representing and manipulating a search
tree.

Existing controlled environments do not isolate backtracking as a learnable
skill.
Grid-world environments such as BabyAI and visual reasoning datasets such as
ARC conflate spatial navigation with logical search, making it impossible to
attribute failures specifically to backtracking.
Formal theorem provers (Lean, Coq) do require systematic backtracking, but
their tactic languages are complex, syntactically overlap with pretraining
corpora, and are tightly coupled to type-theoretic foundations that introduce
confounds for language model experiments.
Tree-of-Thoughts \citep{yao2024tree} and Monte-Carlo Tree Search methods for
LLMs impose a search strategy at inference time but do not train models to
represent or execute the search tree as a syntactic object.

We introduce \gore{} (\textbf{G}raph-\textbf{O}rchestrated \textbf{R}easoning
\textbf{E}nvironment), a minimal formal language designed from the ground up so
that every non-trivial program requires explicit backtracking to evaluate.
\gore{} has five primitives: \textsc{Step} (unification / result binding),
\textsc{Fork} (non-deterministic branching), \textsc{Cut} (labelled branch
pruning), \letstmt{} (deterministic functional computation), and \callstmt{}
(explicit external effect).
The language uses novel syntax with zero overlap with natural-language
pretraining corpora, eliminating leakage as a confound.
A key structural observation underlies our design: \gore{} primitives are
\emph{structurally isomorphic} to standard Lean\,/\,Coq tactic combinators
(\texttt{cases}$\to$\textsc{Fork}, \texttt{exact}$\to$\textsc{Step},
\texttt{contradiction}$\to$\textsc{Cut}, \texttt{apply}$\to$\callstmt{},
\texttt{have}$\to$\letstmt{}), making theorem-proving corpora a natural
pretraining domain.

The \gore{} interpreter serves as a \emph{Process Reward Model} (PRM) without
any additional training.
At every step of execution, the interpreter verifies whether the predicted
trace node (type, environment binding, unification result) is correct and
returns a dense scalar reward.
This stands in sharp contrast to prior PRM work \citep{lightman2023lets}, which
requires either human annotation of intermediate steps or a separately trained
critic model.
The interpreter \emph{is} the PRM: step-level verification is free by
construction.

\gore{}'s \callstmt{} primitive and recursive clause structure formalize the
same recursive task-decomposition pattern that underlies program-aided reasoning
\citep{gao2023pal}, structured query decomposition (Text-to-SQL via CTE trees,
\citealt{yu2018spider}), and the broader ``Decomposer'' paradigm.
In all these settings a complex task is broken into sub-tasks that are
dispatched to tools or sub-agents; \gore{} makes this structure explicit and
verifiable, enabling dense reward at every node of the decomposition tree.

Our contributions are as follows.
\begin{enumerate}
  \item \textbf{\gore{} language.}
    A minimal five-primitive formal language with full operational semantics,
    a reference interpreter, and a parameterised curriculum generator.
  \item \textbf{Structural isomorphism.}
    A formal correspondence between \gore{} primitives and Lean\,/\,Coq tactic
    combinators, enabling warm-start pretraining on existing proof corpora
    (LeanDojo \citep{yang2023leandojo}, CoqGym \citep{yang2019coqgym}).
  \item \textbf{Interpreter-as-PRM.}
    Dense step-level reward from the \gore{} interpreter without human
    annotation or a separate reward model, providing the cleanest possible
    ablation of dense vs.\ sparse reward in an RL-for-reasoning setting.
  \item \textbf{Three-stage training pipeline.}
    Formal proof pretraining $\to$ \gore{} curriculum SFT $\to$ Dr.\,GRPO with
    dense interpreter reward; ablations with GRPO+sparse and DAPO.
\end{enumerate}

The remainder of the paper is organised as follows.
Section~\ref{sec:related} situates \gore{} within the literature.
Section~\ref{sec:methodology} defines the language, interpreter-as-PRM, and
training pipeline.
Section~\ref{sec:experiments} describes the four experimental conditions.
Section~\ref{sec:results} reports results.
Section~\ref{sec:conclusion} concludes with limitations and future work.
